\bookchapter{제45장: 품질 보증 — AI가 쓴 코드를 믿을 수 있게 만들기}

\chapterintro{AI가 쓴 코드에는 한 가지 특징이 있어요: 대개 “맞아 보인다”는 것 — 변수명은 규범적이고 주석은 충실하며 논리는 깔끔해요. 하지만 “맞아 보임”은 “진짜 맞음”이 아니에요. AI 코드에 숨은 문제는 손으로 쓴 코드보다 찾기 어려울 때가 많아요. 표면 품질이 너무 좋아 경계심을 늦추기 쉽기 때문이에요. 이번 장은 AI 코드를 위한 효과적인 품질 보증 체계를 세우는 법을 가르쳐요. — 테스트, 리뷰, 버전 관리라는 세 개의 방어선으로 “맞아 보임”을 “진짜 맞음”으로 바꾸는 거예요.}

\booksection{테스트 주도: AI에게 테스트 시키기}

테스트는 AI 코딩 품질 보증의 첫 번째 방어선이에요. 좋은 소식은: AI에게 테스트를 시키는 편이 기능 코드를 시키는 것보다 쉽고 믿을 만하다는 점이에요. — 테스트의 기대 동작은 명확해서 “맞아 보이지만 실은 틀린” 코드가 나오기 어렵기 때문이에요.

테스트 주도 AI 코딩 흐름은 이래요: 먼저 AI에게 기능 코드를 쓰게 하고, 기본 기능 검증을 통과하면 즉시 테스트 케이스를 쓰게 해요. 테스트 케이스는 정상 경로, 경계 상황, 예외 상황을 다뤄야 해요. 그다음 AI가 테스트 피드백에 따라 기능 코드를 고치게 하죠. 이 흐름의 장점은: 테스트를 한번 만들어 두면 지속적인 검증 도구가 된다는 것 — 이후 어떤 수정이든 테스트를 실행해서 회귀 문제가 없음을 확인할 수 있어요.

\begin{lstlisting}
calculate_discount 함수에 대한 단위 테스트를 작성해 주세요.
테스트 케이스는 다음을 포괄해야 해요:
1. 정상 경로:
   - 100위안 상품 20% 할인 -> 80
   - 200위안 상품 50% 할인 -> 100
2. 경계 사례:
   - 0위안 상품 20% 할인 -> 0
   - 100위안 상품 0% 할인 -> 100(할인 없음)
   - 100위안 상품 10% 할인 -> 100(정가)
3. 예외 사례:
   - 음수 가격 -> 오류를 던져야 함
   - 할인율 > 1 -> 오류를 던져야 함
   - 숫자가 아닌 입력 -> 오류를 던져야 해요
Vitest 프레임워크를 사용해 완전한 테스트 파일을 제시해 주세요.
\end{lstlisting}

이 지시의 특징에 주목하세요: 세 종류의 테스트 케이스(정상, 경계, 예외)를 명확히 나열해 AI가 “골라서 테스트하기”로 재량을 발휘할 여지를 주지 않아요. AI는 가끔 정상 경로 테스트만 써서 “얼렁뚱땅 넘어가려” 하는데, 모든 범주를 명확히 나열하면 이 문제를 피할 수 있어요.

\booksection{코드 리뷰 노하우}

테스트 커버리지가 있어도 AI가 만들어 낸 코드는 사람이 직접 검토해야 해요. 코드 리뷰는 코드를 한 줄씩 통째로 읽는 게 아니에요——그건 너무 시간이 들고 여러분의 코딩 실력에도 크게 의존하거든요. 효과적인 AI 코드 리뷰는 다음의 핵심 항목에 주목해야 해요:

\begin{center}
\begin{tabularx}{\linewidth}{|X|X|X|}
\hline
\textbf{검토 항목} & \textbf{확인할 것} & \textbf{흔한 문제} \tabularnewline
\hline
보안 & 입력값 검증, 시크릿 키 처리, SQL 인젝션 방지 & AI는 입력값 검증을 잊거나 시크릿 키를 코드에 하드코딩하곤 해요 \tabularnewline
오류 처리 & 예외를 잡아 내는지, 오류 메시지가 적절한지 & AI는 빈 catch 블록이나 console.log만 쓰는 걸 좋아해요 \tabularnewline
성능 & N+1 쿼리, 무한 루프 같은 뚜렷한 성능 문제가 있는지 & AI는 가끔 O(n²)짜리 알고리즘을 만들어 내기도 해요 \tabularnewline
유지보수성 & 코드 구조가 명확한지, 이름이 의미 있게 지어졌는지 & AI는 가끔 지나치게 추상적인 코드를 만들어 내요 \tabularnewline
의존성 관리 & 가져온 의존성이 꼭 필요한지, 버전이 명확한지 & AI는 불필요한 라이브러리를 가져오는 경향이 있어요 \tabularnewline
\hline
\end{tabularx}
\end{center}

\begin{center}\small\color{gray}AI 코드 리뷰의 핵심 항목\end{center}

“이 코드의 보안, 오류 처리, 성능을 검토하고 잠재적인 문제를 짚어 줘.”라고 지시하면 돼요. AI의 자체 검토 능력은 계속 좋아지고 있어서, 사람이 하는 검토를 완전히 대신하지는 못하지만 뻔히 보이는 문제 대부분을 걸러 주는 데 도움이 돼요.

\booksection{버전 관리: AI 코딩에서 Git의 역할}

AI 코딩에서 Git의 역할은 전통적인 프로그래밍 때보다 더 중요해요. AI가 수정할 때마다 예측 못 할 변화가 생길 수 있으니 믿을 수 있는 '후회 해독제'가 필요한데——버전 관리가 바로 그 후회 해독제예요.

AI 코딩에서 Git의 모범 사례는 '자주 커밋하고, 작게 나눠 빠르게 나아가기'예요. 검증할 수 있는 기능 하나를 마칠 때마다 커밋하세요. 그러면 AI의 수정에 문제가 생겼을 때 모든 작업을 엎지 않고도 바로 이전에 동작하던 상태로 정확히 되돌릴 수 있어요.

또 하나 유용한 팁은 '브랜치 격리'예요. 결과가 어떨지 모르는 AI 수정(예: 큰폭의 리팩터링)은 새 브랜치를 만들어서 진행하세요. 결과가 마음에 들면 기본 브랜치로 병합하고, 마음에 들지 않으면 브랜치만 지우면 되니까 기본 브랜치에 영향을 주지 않아요. 이게 기본 브랜치에서 여러 번 고치는 것보다 훨씬 안전해요.

\begin{enumerate}
  \item AI에게 수정을 맡기기 전에는 작업 공간이 깨끗한지(커밋하지 않은 변경이 없는지) 확인하세요. 그래야 AI의 수정이 마음에 들지 않을 때 git checkout . 으로 빠르게 되돌릴 수 있어요.
  \item AI의 수정이 끝나고 검증을 통과하면 바로 커밋하세요. 커밋 메시지에는 이번 수정에서 무엇을 했는지 적는데, 예컨대 "feat: 사용자 회원가입 기능 추가"처럼요.
  \item 큰폭의 리팩터링은 먼저 브랜치부터 만드세요: git checkout -b refactor/auth-system. 브랜치에서 수정을 마치고, 테스트를 통과한 뒤에 병합하세요.
  \item diff를 확인하는 습관을 들이세요. AI가 수정할 때마다 git diff로 무엇을 바꿨는지 구체적으로 살펴 보는 거죠. 문제를 발견하는 데도 도움이 되고, 코드를 배우는 데도 도움이 돼요.
\end{enumerate}

\booksection{자주 겪는 함정과 피하는 법}

AI 코딩에는 몇 가지 되풀이되는 함정이 있는데, 초보자는 거의 전부 밟게 돼요. 미리 알아 두면 헤매는 길을 크게 줄일 수 있어요.

\booksubsection{함정 하나: 맹신의 함정}

“AI가 쓴 코드가 꽤 전문가처럼 보이니까 분명 괜찮겠지?”——이게 가장 위험한 함정이에요. AI가 만들어 낸 코드는 겉보기에 아주 깔끔해요——이름 짓기 규칙도 지켰고, 주석도 빠짐없고, 구조도 명확하고요. 이런 '겉보기 전문성'이 사람의 경계심을 늦춰요. 하지만 겉이 깔끔하다고 논리가 맞는 건 아니에요. 해독제: '그래 보이는가'로 판단하지 말고 항상 테스트로 검증하세요.

\booksubsection{함정 둘: 비대해짐의 함정}

“AI가 코드를 대신 짜 주는데, 차라리 기능을 풍성하게 만들면 어때?”——이러면 코드가 비대해져요. AI는 '더 주기'를 아주 잘해요. 로그인 페이지 하나만 만들라고 해도 사용자 관리 시스템 통째로 만들어 주곤 하죠. 불필요한 기능은 복잡도를 늘릴 뿐 아니라 버그가 날 확률도 높여요. 해독제: Simplicity First 원칙을 철저히 지키고, AI에게 “이 기능만 만들고, 다른 건 덧붙이지 마”라고 분명히 말하세요.

\booksubsection{함정 셋: 이해의 함정}

“AI가 알아서 내 뜻을 이해하겠지?”——꼭 그렇지는 않아요. AI의 이해력이 아무리 뛰어나도 모호한 요구 앞에서는 여전히 '추측'해요. 추측이 틀리면 나온 코드도 틀린 거고요. 더 골치 아픈 건, 틀리게 추측한 코드도 대개 '그럴듯해 보여서' 이해했다고 착각하게 만든다는 점이에요. 해독제: 지시에 구체적인 세부사항을 충분히 담아서 AI가 '추측'할 여지를 줄이세요.

\booksubsection{함정 넷: 맥락의 함정}

“예전 대화에서 말했으니 AI가 기억하고 있겠지?”——아마 기억하지 못해요. 대화가 길어지거나 새 세션을 열면 AI는 앞의 맥락을 잃어버려요. 중요한 제약 조건이 프로젝트 기억 파일에 적혀 있지 않으면 뒤이은 수정에서 무시되기 쉬워요. 해독제: 핵심 제약은 프로젝트 기억 파일에 적어 두고, 대화 기억에 의존하지 마세요.

\begin{tcbwarn}{네 가지 함정 요약}
맹신의 함정——'그래 보인다'고 믿지 마세요. 비대해짐의 함정——AI는 '더 주기'를 좋아해요. 이해의 함정——모호한 요구는 '추측'해요. 맥락의 함정——대화 기억은 '잊어버려요'. 해독제는 차례로 테스트 검증, 불필요한 것 잘라내기, 세부사항 충분히 주기, 기억 파일에 적기예요.

\end{tcbwarn}

\begin{tcbpractice}{따라 해 봐요: 프로젝트에 3중 방어선 세우기}
프로젝트 1(도구 사이트)을 이미 시작했다면: (1) AI에게 핵심 함수의 테스트 케이스 3종을 작성하게 하기; (2) 다섯 가지 항목(보안/오류 처리/성능/유지보수성/의존성)으로 코드를 한 번 검토하기; (3) git으로 프로젝트를 초기화하고 첫 커밋 하기. 이 3중 방어선이 갖춰져야 프로젝트가 진짜 여러분의 것이 돼요.

\end{tcbpractice}

\begin{tcbsummary}{이번 장 정리}
\begin{itemize}
  \item AI 코드가 '그래 보인다' $\neq$ '진짜 맞다'——겉 품질이 너무 좋으면 오히려 경계심이 느슨해지기 쉬워요
  \item AI에게 테스트 작성 맡기기: 정상/경계/예외 세 종류의 케이스를 분명히 나열해서 '골라 내는 테스트'를 막기
  \item 검토 다섯 항목: 보안, 오류 처리, 성능, 유지보수성, 의존성 관리
  \item Git은 '후회 해독제': 자주 커밋, 작게 나눠 빠르게, 브랜치 격리, diff 자주 보기
  \item 네 가지 함정: 맹신(테스트 검증), 비대해짐(불필요한 것 잘라내기), 이해(세부사항 충분히), 맥락(기억 파일에 적기)
\end{itemize}

\end{tcbsummary}
